% =====================================================================
% Paper: Sentence-Level Humor Detection on Stand-Up Comedy:
%        A 5x3 Repeated Cross-Validation Study
% =====================================================================
% Authors: Subhajit Das (Independent Researcher)
% Date:    2026-10-08
% Status:  Draft for arXiv preprint (cs.CL) / INTERSPEECH 2026 / EMNLP 2026
% Repo:    https://github.com/Das-rebel/HaHaScore  (commit 65a65cc)
% Contact: sdas22@gmail.com
% =====================================================================

\documentclass[11pt, letterpaper]{article}

\usepackage[utf8]{inputenc}
\usepackage[T1]{fontenc}
\usepackage{amsmath, amssymb}
\usepackage{booktabs}
\usepackage{graphicx}
\usepackage{hyperref}
\usepackage{microtype}
\usepackage{multirow}

\title{Sentence-Level Humor Detection on Stand-Up Comedy:\\
       A 5$\times$3 Repeated Cross-Validation Study}

\author{Subhajit Das \\
Independent Researcher \\
\texttt{sdas22@gmail.com}}

\date{\today}

\begin{document}
\maketitle

\begin{abstract}
We study the problem of sentence-level humor detection on stand-up comedy using multimodal fusion of text and audio features. We train and evaluate a series of models on 639 stand-up comedy videos from StandUp4AI, comparing text-only RoBERTa, text+audio bilinear fusion (v5), and a temporal BiGRU (Bridge 4) that incorporates WavLM-base-plus acoustic embeddings with prosodic features. Our central empirical contribution is a 5$\times$3 repeated joke-disjoint cross-validation of the Bridge 4 architecture, which yields an honest Spearman $\rho = 0.9285 \pm 0.0084$ (95\% bootstrap CI $[0.9246, 0.9331]$, $n=15$ measurements) --- substantially above the single-fold 5$\times$ CV figure of 0.8400 ($\Delta = +0.086$) reported in earlier documentation. We further document the methodological lesson that single-fold 5$\times$ CV on 639 stand-up videos can either under-report or over-report model quality, and that the bootstrap-validated 5$\times$3 mean is the only honest number to cite. A series of negative results --- Bridge 5 iterative self-training at AUC 0.97, RoBERTa single-fold 0.32, DeBERTa-v3-small 0.29 --- illustrate the spectrum of fold-luck and architectural choice artefacts that the 5$\times$3 protocol is designed to detect.
\end{abstract}

% =====================================================================
\section{Introduction}
% =====================================================================

Sentence-level humor detection on stand-up comedy is a small but persistent sub-field of speech classification, with applications in content moderation, recommendation, caption timing, and accessibility. The standard approach is multimodal: text (RoBERTa or DeBERTa) + audio (WavLM or wav2vec2) + a fusion head. Reported numbers range from 0.50 (text-only is RANDOM) to 0.85 (multimodal cascade with temporal context).

A common methodological pitfall in this sub-field is reporting single-fold 5$\times$ cross-validation numbers as headline metrics. On 639 stand-up videos with 20 segments each, the within-fold variance of binary humor AUC is approximately 0.03 --- comparable to the entire gap between a text-only baseline (0.50) and a strong multimodal model (0.85). A single-seed 5$\times$ CV at $\rho = 0.85$ is statistically indistinguishable from a single-seed 5$\times$ CV at $\rho = 0.80$.

We apply 5$\times$3 repeated joke-disjoint cross-validation with bootstrap confidence intervals to the Bridge 4 architecture and report the honest mean. Our findings are:
\begin{itemize}
\item \textbf{Bridge 4 honest 5$\times$3 AUC} = $0.9285 \pm 0.0084$ (95\% CI $[0.9246, 0.9331]$, $n=15$). This is \textbf{+0.086} above the single-fold 5$\times$ CV figure of 0.8400 reported in earlier documentation. The single-fold number was \emph{conservative}, not inflated.
\item \textbf{Bridge 5 iterative self-training} at single-fold 5$\times$ CV reports AUC = 0.9685 --- structurally identical to the fold-luck pattern documented in our prior work, and the honest 5$\times$3 mean would be substantially lower.
\item \textbf{Text-only RoBERTa baseline} at single-fold $\rho = 0.32$ (v7), 0.28 (v8), 0.29 (v9) --- all consistent with the text-only RANDOM finding from earlier documentation.
\end{itemize}

The methodological contribution of this paper is the application of 5$\times$3 repeated joke-disjoint CV to a sentence-level humor detection task, with a worked example where the single-fold number was conservative (the opposite of the fold-luck pattern that falsified earlier per-language normalization claims).

% =====================================================================
\section{Setup}
% =====================================================================

\subsection{Data}

The dataset is 639 stand-up comedy videos from StandUp4AI~\cite{standup4ai}, with pseudo-labels generated by an automatic laughter detector applied to 12,780 segments (20 segments per video). Of the 639 videos, 8 are also annotated with per-segment gold labels for the audience-reaction time spans. Per-video pseudo-label confidence is computed as the mean detector output over the 20 segments; videos with confidence $> 0.3$ are treated as confident (84 of 639 in the Bridge 4 baseline, varying by iteration in the self-training variants).

\subsection{Models}

We compare four model families:
\begin{enumerate}
\item \textbf{RoBERTa text-only} (v7--v9): 125M parameters, $\rho \in [0.28, 0.32]$ on single-fold 5$\times$ CV.
\item \textbf{v5 multimodal} (DeBERTa-v3-base + WavLM-base-plus bilinear): val AUC 0.613.
\item \textbf{v10 multimodal cascade}: 5$\times$3 honest AUC 0.6924 humor / 0.5453 gold.
\item \textbf{Bridge 4} (BiGRU 128d, 2-layer, bidirectional over 791d input = 768d WavLM-base-plus + 23d prosody + 4d position): single-fold AUC 0.8400, honest 5$\times$3 AUC 0.9285.
\item \textbf{Bridge 5} (Bridge 4 + iterative self-training): single-fold AUC 0.9685.
\end{enumerate}

The Bridge 4 and Bridge 5 features were extracted once and cached at \texttt{/tmp/bridge4\_features.npz} (44.9 MB) to enable cheap re-evaluation with different validation protocols.

\subsection{Validation protocol}

We report three numbers for each model:
\begin{itemize}
\item \emph{Single-fold 5$\times$ CV} at the seed used in the original documentation: 5 folds at one seed.
\item \emph{Honest 5$\times$3 repeated joke-disjoint CV}: 3 seeds $\times$ 5 folds = 15 measurements. Jokes never appear in both train and test.
\item \emph{95\% bootstrap CI} on the 15 measurements.
\end{itemize}

% =====================================================================
\section{Results}
% =====================================================================

\subsection{Headline: Bridge 4 honest 5$\times$3}

Table~\ref{tab:main} compares the headline numbers.

\begin{table}[h]
\centering
\begin{tabular}{lcc}
\toprule
\textbf{Validation} & \textbf{Bridge 4 AUC} & \textbf{Notes} \\
\midrule
Single-fold 5$\times$ CV (per \texttt{bridge4\_complete}) & 0.8400 & 5 folds at one seed \\
\midrule
5$\times$3 repeated honest mean & \textbf{0.9285} $\pm$ 0.0084 & 3 seeds $\times$ 5 folds = 15 measurements \\
95\% bootstrap CI & $[0.9246,\ 0.9331]$ & Per-fold std 0.0084 \\
\midrule
$\Delta$ (5$\times$3 $-$ single-fold) & $+0.086$ & 5$\times$3 mean is HIGHER \\
\bottomrule
\end{tabular}
\caption{Bridge 4 humor-detection AUC under single-fold vs.\ 5$\times$3 repeated joke-disjoint CV. The 5$\times$3 mean is $0.086$ above the single-fold number --- a \emph{conservative} single-fold in this case, opposite of the fold-luck pattern documented elsewhere.}
\label{tab:main}
\end{table}

The per-seed means (across the 5 folds) are: seed=42 mean 0.9300, seed=142 mean 0.9250, seed=242 mean 0.9305. The inter-seed range is only 0.005, indicating that the 5$\times$3 mean is stable across joke-disjoint randomisations.

\subsection{Comparison to related work}

Table~\ref{tab:compare} compares Bridge 4 to the earlier v5/v7/v8/v9/v10 numbers.

\begin{table}[h]
\centering
\begin{tabular}{lccc}
\toprule
\textbf{Model} & \textbf{Metric} & \textbf{Single-fold} & \textbf{Honest 5$\times$3} \\
\midrule
v7 (RoBERTa text) & $\rho$ & 0.3226 & not run \\
v8 (RoBERTa text, cosine LR) & $\rho$ & 0.2820 & not run \\
v9 (DeBERTa-v3-small text) & $\rho$ & 0.2892 & not run \\
\midrule
v5 (DeBERTa + WavLM bilinear) & AUC & 0.613 (Val) & not run \\
v10 (multimodal cascade) & AUC & 0.6924 humor / 0.5453 gold & 0.6924 / 0.5453 \\
\midrule
\textbf{Bridge 4 (BiGRU WavLM + prosody)} & \textbf{AUC} & \textbf{0.8400} & \textbf{0.9285 $\pm$ 0.0084} \\
\bottomrule
\end{tabular}
\caption{Comparison of humor-detection models on StandUp4AI. Bridge 4's honest 5$\times$3 result (0.9285) is the only number that has been bootstrap-validated.}
\label{tab:compare}
\end{table}

\subsection{Why Bridge 4's 5$\times$3 mean exceeds the single-fold number}

In our prior work on a per-language normalization claim, a single-fold measurement of $+0.163$ was found to be the result of fold-luck --- the 5$\times$3 repeated CV showed a mean of $-0.127$ (a swing of $0.290$). The lesson is that single-fold numbers on imbalanced, joke-disjoint data can be inflated by fold-luck.

For Bridge 4, the direction is opposite. The single-fold 5$\times$ CV number of 0.8400 under-represents the model's performance by $0.086$. The 15-measurement distribution is tight (per-fold std $0.0084$), so the bootstrap CI $[0.9246, 0.9331]$ does not overlap the single-fold number.

This does not contradict the prior fold-luck finding --- it is a separate, more-favourable case of the same methodological lesson. The honest number is the 5$\times$3 mean in both directions; the single-fold number can be either above or below the truth.

% =====================================================================
\section{Discussion}
% =====================================================================

\subsection{Bridge 5 iterative self-training at AUC 0.97 is fold-lucky}

The Bridge 5 self-training variant reports single-fold AUC 0.9685 across 133 confident files. Per the methodological lesson, this is the same pattern as the $+0.163$ per-language norm claim: a single-seed 5$\times$ CV measurement can be inflated by $0.10$ or more. The honest 5$\times$3 mean for Bridge 5 has not been computed in this paper; doing so is the natural next step. Per the predicted outcome from the structural argument, the 5$\times$3 mean for Bridge 5 is expected to be substantially lower than 0.97, possibly in the 0.85--0.90 range.

The lesson: \textbf{any single-fold AUC $> 0.95$ on 133--639 stand-up videos should be treated as fold-luck until a 5$\times$3 honest measurement is available}. We apply this lesson to ourselves: we do not report Bridge 5's 0.9685 as a headline number.

\subsection{Implications for the HaHaScore product vision}

The product vision is a 3-product platform per the recovered PRD v6: (1) Group Laughter Predictor (word-level discourse positioning, IoU-F1 metric), (2) Individual Laughter + Sarcasm (deferred, needs diarization), (3) Sarcasm-Aware Content Scorer. The Bridge 4 honest 5$\times$3 = 0.9285 establishes the sentence-level humor detection baseline against which future P1 (word-level discourse positioning) extensions can be measured.

\subsection{Limitations}

\begin{enumerate}
\item \textbf{Single-label binary}: The Bridge 4 AUC is binary humor/not-humor. P1 requires word-level IoU-F1, which is a different metric.
\item \textbf{English only}: StandUp4AI is English-only standup. Generalization to other languages and comedy styles is out of scope.
\item \textbf{8 of 639 videos have gold labels}: The honest 5$\times$3 is over a pseudo-label task; the 8 gold-labeled videos are used for evaluation only. Direct validation on the 8 gold-labeled subset is left to future work.
\item \textbf{Single architecture}: We evaluated only Bridge 4 with 5$\times$3 honest. Other models (v5, v10) would also benefit from 5$\times$3 honest evaluation.
\end{enumerate}

% =====================================================================
\section{Conclusion}
% =====================================================================

We apply 5$\times$3 repeated joke-disjoint cross-validation with bootstrap confidence intervals to the Bridge 4 sentence-level humor detection architecture on StandUp4AI, and report the honest mean AUC = $0.9285 \pm 0.0084$ (95\% CI $[0.9246, 0.9331]$, $n=15$ measurements). The single-fold 5$\times$ CV figure of 0.8400 is conservative by $0.086$, the opposite direction of the fold-luck pattern that falsified earlier per-language normalization claims. We recommend that all future HaHaScore headline numbers be reported via the 5$\times$3 honest protocol with bootstrap CIs, regardless of whether the single-fold direction is inflation or deflation.

% =====================================================================
% Reproducibility
% =====================================================================

The repository at \texttt{github.com/Das-rebel/HaHaScore} (commit \texttt{65a65cc}) contains:
\begin{itemize}
\item \texttt{eval\_bridge4\_5x3.py}: the 5$\times$3 honest CV evaluation script.
\item \texttt{experiments/bridge4\_5x3/result.json}: the 15 measurements and bootstrap CI.
\item \texttt{experiments/bridge4\_5x3/RESULT.md}: human-readable result.
\item Cached Bridge 4 features at \texttt{/Users/Subho/tmp/bridge4\_features.npz} (44.9 MB).
\end{itemize}

% =====================================================================
% References
% =====================================================================
\begin{thebibliography}{9}

\bibitem{standup4ai}
StandUp4AI: A Multi-Cultural Humor Detection Dataset.
\newblock \textit{arXiv preprint}, 2024.

\bibitem{roberta}
Y. Liu, et al.
\newblock RoBERTa: A Robustly Optimized BERT Pretraining Approach.
\newblock \textit{arXiv preprint}, 2019.

\bibitem{wavlm}
S. Chen, et al.
\newblock WavLM: Large-Scale Self-Supervised Pre-Training for Full Stack Speech Processing.
\newblock \textit{arXiv preprint}, 2022.

\bibitem{deberta}
P. He, et al.
\newblock DeBERTaV3: Small Models Tackling Major NLP Tasks.
\newblock \textit{arXiv preprint}, 2023.

\bibitem{hahascore-v10}
S. Das.
\newblock HaHaScore v10: Multimodal Cascade with Speaker-Disjoint CV.
\newblock HaHaScore repo, \texttt{models/v10\_cascade}, 2026.

\bibitem{hahascore-v5}
S. Das.
\newblock HaHaScore v5: DeBERTa-v3-base + WavLM-base-plus Bilinear Fusion.
\newblock HaHaScore repo, 5-fold CV AUC 0.632 $\pm$ 0.007, 2026.

\bibitem{hahascore-bridge4}
S. Das.
\newblock Bridge 4: BiGRU over 791d Multimodal Features.
\newblock HaHaScore repo, \texttt{models/bridge4}, single-fold AUC 0.8400 $\pm$ 0.027, 2026.

\bibitem{hahascore-falsification}
S. Das.
\newblock Evaluating Speech-Humor Models Without Identity Leakage: A Self-Falsification Case Study.
\newblock \textit{arXiv preprint}, 2026. Documents the $+0.163 \to -0.127$ per-language normalization falsification that motivates the 5$\times$3 honest CV protocol.

\bibitem{hahascore-v1}
S. Das.
\newblock Continuous Humor-Strength Prediction from Text: A 5$\times$3 Repeated Cross-Validation Baseline on Jester.
\newblock \textit{arXiv preprint}, 2026. Companion paper; documents the Jester text-only continuous regression task (v1 honest 5$\times$3 $\rho = 0.2292$).

\end{thebibliography}

\end{document}
